\clearpage
\section{Triangular bunch in a conducting circular pipe}
\label{sec:triangle-pipe}
This control reproduces the longitudinal-force component of the straight
round-pipe triangular-profile benchmark used in ORBIT~\cite{orbit2006,pyorbit3triangle}. It tests a local
longitudinal field in a prescribed source; it is not a curved-field or tracking
validation. The source has fixed circular radius $a=5$ mm, full length
$L=10$ m, $N=10^{14}$ protons and kinetic energy $0.4$ GeV. The grounded
pipe radius is $b=20$ mm. All comparisons use a single equivalent frozen
kick of length $\ell=1$ m.

With $s=0$ at the bunch centre, the line density and long-beam field are
\begin{align}
 \lambda(s)&=\frac{2Q}{L}\left(1-\frac{2|s|}{L}\right),
   &&|s|<L/2,\\
 E_s(r,s)&=-\frac{\lambda'(s)}{4\pi\eps\gamma^2}
 \left[1+2\ln\frac ba-\frac{r^2}{a^2}\right],
   &&0\le r\le a.
 \label{eq:triangle-pipe}
\end{align}
The density vanishes outside its support. The bracket is $g_{\rm axis}
=1+2\ln(b/a)$ on axis and has uniform-disk mean
$\overline g=1/2+2\ln(b/a)$~\cite{baartman}. The local formula is compared
only at $1\le|s|\le4$ m, excluding the cusp and ends. A proton's equivalent
energy kick is $\Delta W[\mathrm{eV}]=\ell E_s[\mathrm{V/m}]$.
The common physical reference gives $E_*=|E_s(0,s)|=10.681199$ kV/m on
either linear branch. No coefficient is fitted.

\reportfigure{triangle-pipe-fields}{Frozen on-axis field and signed difference
from \cref{eq:triangle-pipe}. Black dashed: physical analytic solution;
symbols: native results. Gray regions are excluded from acceptance;
residuals use the fixed scale $E_*$ and the gray band marks $\pm1\%$.
The cusp transition is displayed without an accuracy claim.}
{fig:triangle-pipe-fields}
\tabfile{triangle-results}
Both columns are maximum pointwise relative differences from the physical
analytic field: 62 on-axis targets and 82 radial targets at $s=\pm2.5$ m.
All accepted reference values are nonzero. PyORBIT's sliced-potential and
Baartman variants meet 1\%; native OPALX CIRCULAR does not.

\clearpage
\subsection{Numerical controls and the remaining model difference}
All numerical changes quoted below use the 144 accepted targets.
\reportfigure{triangle-pipe-controls}{Left: the physical analytic radial
dependence and the selected native runs in \cref{fig:triangle-pipe-fields}.
Right: physical on-axis errors under selected mesh changes. PyORBIT varies
the transverse grid at fixed $N_z=128$ and the same 16-ring source; OPALX
varies the coupled mesh from $128^3$ to $256^3$ with the same finer source.
The primary PyORBIT slice result on the left uses 32 rings. The dotted line
marks 1\%; these controls are not a joint asymptotic extrapolation.}
{fig:triangle-pipe-controls}

PyORBIT's sliced-potential variant imposes a grounded circular boundary and
explicitly enables longitudinal kicks. Its last transverse mesh change is
0.3384\% of $E_*$ on accepted targets; doubling radial and angular source sampling changes the
field by 0.2838\%. The remaining on-axis bias, $-0.2861\%$, agrees with an
independent discrete-ring potential prediction, $-0.2860\%$. The Baartman
variant uses the radial $g(r)$ in \cref{eq:triangle-pipe}. Its largest
accepted physical error is 0.00660\%.

OPALX uses \code{PIPEMODE=CIRCULAR}, \code{BEAMR=0.005} m and
\code{PIPESIZEX=PIPESIZEY=0.020} m. These sizes are full numerical widths
$W$ of the open transverse mesh, while the longitudinal formula uses
\begin{equation}
 g_{\rm OPALX}=0.67+2\ln(W/a)=3.442589,
 \qquad g_{\rm axis}=3.772589.
 \label{eq:triangle-native}
\end{equation}
Thus $W=b$ matches the effective pipe radius in this formula, although its
open numerical grid spans $\pm10$ mm; no transverse conducting wall is solved.
For this centred round source the local transverse image field vanishes.
The coefficient 0.67 belongs to an equivalent parabolic-bunch moment
model~\cite{oeftiger}, rather than this pointwise fixed-radius disk reference.
OPALX also omits the radial dependence of $E_s$. These model differences
predict an on-axis bias of $-8.7473\%$ before numerical errors. The refined
run follows its native formula within 0.161\%; the matched mesh change is
0.1253\% of $E_*$ on accepted targets. Increasing $W$ to 40 mm at matched transverse spacing
changes the native $g$ to 4.828883. This changes the longitudinal model,
so it is not a padding test at fixed physical pipe radius.

All six OPALX source/probe pairs and 16 PyORBIT configurations pass independent
input, charge, position and field-conversion checks. OPALX uses a
$10^{-20}$ s field readout and subtracts a matched probe-only solve;
maximum probe motion is $2.2\times10^{-12}$ m. PyORBIT uses zero-charge
probes and fixed-position native kicks. One-/two-rank PyORBIT fields agree
within $1.5\times10^{-12}\%$ of $E_*$; FFT2D5 remains single-rank.
An unresolved $N_z=256$ source-sampling case is retained in the archived data;
increasing longitudinal samples removes its alias. No production solver,
coefficient or physical tolerance was changed. These results concern the
longitudinal field only, with no tune-shift, full-drift or curved-field claim.
